%
% einleitung.tex
%
% (c) 2026 Nathanael Fässler, OST Ostschweizer Fachhochschule
%
% !TEX root = ../../buch.tex
% !TEX encoding = UTF-8
%
\section{Lineare Gleichungssysteme
\label{buchberger:section:teil0}}
\kopfrechts{Lineare Gleichungssysteme}
Viele Elemente des Gauss-Algorithmus werden in einer ähnlichen Form im Abschnitt zum Buchberger-Algorithmus~\ref{buchberger:section:buchberger} wiederkehren.
Daher ist es empfehlenswert die Grundideen des Lösens eines linearen Gleichungssystems zu repetieren.

Das Gleichungssystem
\begin{equation}
  \left\{ % macht auf linker Seite die grosse geschweifte Klammer (darum das \{)
  \renewcommand{\arraycolsep}{2pt} % überschreibt Spaltenabstand in Tabelle auf 2pt, sonst wären sie zu gross
  \begin{array}{rcrcr} % Array == Tabelle. rcrcr definiert die 5 Spalten. r = rechtsbündig, c = centered. Z.B. x und 3x zentrieren sich nach rechts zum +
    x &+& 2y &=& 3 \\
    3x &-& y &=& 15
  \end{array}
  \right. % macht auf rechter Seite keine Klammer, darum der Punkt.
  \label{buchberger:lineares-gleichungssystem-1}
\end{equation}

lässt sich bekanntlich mit dem Gauss-Algorithmus lösen.
Das Vorgehen ist wie folgt:
\begin{enumerate}
\item
Die Gleichungen werden in eine passende Form gebracht:

Alle Terme mit einer Variablen werden auf die linke Seite verschoben, während auf der rechten Seite der konstante Wert steht.
Die Reihenfolge der Variablen muss einer festgelegten Ordnung folgen, damit der Algorithmus deterministisch ist.
Das Gleichungssystem~\eqref{buchberger:lineares-gleichungssystem-1} ist bereits in dieser Form, wobei die Ordnung dadurch gegeben ist, dass $x$ vor $y$ steht.
\item
Mit iterativen Gleichungsumformungen, die hier nicht weiter relevant sind, wird das Gleichungssystem in eine treppenartige Form gebracht.
\begin{equation}
  \left\{
  \renewcommand{\arraycolsep}{2pt}
  \begin{array}{rcrcr}
  x &+& 2y &=& 3 \\
     &&  y &=& -\frac{6}{7}
  \end{array}
  \right.
  \label{buchberger:lineares-gleichungssystem-2}
\end{equation}
\item
Aus dem Gleichungssystem~\eqref{buchberger:lineares-gleichungssystem-2} lässt sich direkt die Lösung für $y$ ablesen, welche $-\frac{6}{7}$ beträgt.
Diese Lösung für $y$ lässt sich in die obere Gleichung einsetzen.
Dabei ensteht eine Gleichung, die nur noch $x$ enthält. Das Auflösen nach $x$ ist trivial.
\end{enumerate}

Der Algorithmus lässt sich auf beliebig grosse Gleichungssysteme anwenden:

Bei einem Gleichungssystem mit den Variablen $x_{0}, x_{1}, \dots, x_{n}$ erhält man nach der Vorwärtsreduktionsphase des Gauss-Algorithmus eine Zeilenstufenform,
wobei in der untersten Gleichung nur noch die Variable $x_{n}$ vorkommt und deren Lösung sich direkt ablesen lässt.
Die nächsthöhere Gleichung lässt sich dann iterativ lösen, indem man immer die Lösungen der bereits fixierten Variablen einsetzt.\footnote{Diese letzte Phase des Gauss-Algorithmus nennt sich das \emph{Rückwärtseinsetzen}.}

Bemerkenswert ist, dass ein zunächst schwer lösbares Gleichungssystem in ein äquivalentes\footnote{Zwei Gleichungssysteme sind äquivalent, wenn sie exakt dieselbe Lösungsmenge besitzen.} Gleichungssystem umgewandelt wird, welches einfacher lösbar ist.